% \begin{frame}{Conclusion}
% \framesubtitle{Main contribution and implications}

% \begin{block}{Main finding}
% Formative LLM measurement is useful, but construct-dependent: it works best when the quality construct has empirically separable dimensions.
% \end{block}

% \vspace{0.3cm}

% \begin{itemize}
%     \item LLM evaluation should be treated as measurement, not only prediction.
%     \item Human ratings and behavioral outcomes provide validity evidence, not true scores.
%     \item Fine-grained criteria can improve both validity evidence and diagnostic interpretation.
% \end{itemize}

% \end{frame}


\begin{frame}{Conclusion}
\framesubtitle{Moving from holistic prediction to meaningful measurement}

\begin{block}{The Core Takeaway}
Formative LLM measurement provides stronger validity evidence than holistic baselines, but its utility is construct-dependent: it works best when the dimensions are empirically separable.
\end{block}

\vspace{0.1cm}

\textbf{Broader Implications for the Field:}
\vspace{0.05cm}
\begin{itemize}
    \setlength{\itemsep}{0.3cm} % Adds breathing room between bullet points
    \item \textbf{Epistemic shift:} LLM evaluation should be treated as formal measurement, not merely text prediction.
    \item \textbf{Validity framing:} Human ratings and behavioral outcomes provide validity evidence, not infallible ``true scores.''
    \item \textbf{Practical design:} Explicit, fine-grained criteria improve both statistical validity and diagnostic transparency.
\end{itemize}

\end{frame}